% 第 36 章。数字与教学模型：checks/check_ch36.py；图：figsrc/fig_ch36.py。
% 史料：brettonwoods_facts（取数脚本 data/fetch/brettonwoods36.py）。
\chapter{布雷顿森林体系的运行}\label{ch:brettonwoods}

\begin{chapterguide}
\guidequestion{美元为什么既是体系的支柱，又是它的隐患？}
\begin{guidegoals}
\item 描述固定但可调整的汇率、资本管制与国际货币基金组织的角色；
\item 解释 1958 年西欧恢复经常账户可兑换以后的变化；
\item 推导特里芬难题（1960）；
\item 分析伦敦黄金总库与 1960 年代的美元危机。
\end{guidegoals}
\guideprereq{\chref{ch:warfinance}\enspace\nameref{ch:warfinance}；\chref{ch:external}\enspace\nameref{ch:external}}
\end{chapterguide}

1965 年 2 月 4 日，法国总统戴高乐（Charles de Gaulle，1890—1970）在记者招待会上回答提问，呼吁世界货币体系逐步回到金本位\cite{frus36,loc36}。照美国国务院外交文件的概括，他主张负有国际金融责任的国家考虑只用黄金清偿国际收支的债务\cite{frus36}。他说，黄金是“唯一无可置疑的、不带任何一国印记的货币基础”，所以是进行国际金融交易最公平的标准\cite{frus36}。
% 据: fr36_date, loc_degaulle, fr36_debts, fr36_gold, bo_campaign

同一天，美国财政部发表声明反对。声明说，那样的货币体系 1931 年已经“崩溃”，“无力为世界贸易的巨大增长提供资金”\cite{frus36}。
% 据: fr36_treasury

这一年年初，法国已宣布对外的盈余今后一律收取黄金\cite{bis36}。在美英两国逆差的背景下，许多人由此想到金价可能上涨，“随后是大量投机性的买金”\cite{bis36}。1965 年，美国向西欧净售出近 13 亿美元的黄金，其中卖给法国 8.8 亿美元\cite{bis36}。美国的黄金存量这一年减少了约 17 亿美元，西欧的黄金储备增加了约 19 亿美元\cite{bis36}。
% 据: b66_france, b66_speculation, b66_table, b66_sold, b66_decline

法国能这样做，是因为体系本来就这样安排。美国是唯一按黄金钉住本国货币、买卖黄金的国家，其他国家都对美元定下平价，也就是官方规定的本币与美元的比价，并在平价上下 1\% 以内干预汇率\cite{bordo1993}。1946 年起，金价是每盎司 35 美元\cite{bordo1993}。盎司是称量黄金的重量单位（\chref{ch:goldstandard}）。外国的货币当局攒下美元，可以拿到美国财政部换成黄金\cite{bordo1993}。美元为什么既是体系的支柱，又是它的隐患？
% 据: bo_pegged, bt_us, bo_pool40, bo_convert

\section{当时的世界}\label{sec:brettonwoods-world}

英美两国战时的方案，在\secref{sec:warfinance-plans}；1944 年议定的协定，包括可调整的平价、份额、稀缺货币条款与资本管制，在\secref{sec:warfinance-bw}。1946 年 12 月，32 国宣布了官方平价\cite{bordo1993}。大战结束时，美国握有世界货币黄金的三分之二\cite{bordo1993}。从宣布平价到主要工业国恢复经常账户交易的可兑换，用了 12 年\cite{bordo1993}。经常账户的可兑换，是说为进出口、运费、旅行这类经常交易，本币可以自由换成外币；为投资与借贷换汇属于资本账户，协定允许各国管制（\secref{sec:warfinance-bw}）。
% 据: bo_parvalues, bo_twothirds

缺的是美元。欧洲经济合作组织各国的国际收支逆差，因 1946—1947 年冬天的严寒而加重，1947 年高达 90 亿美元\cite{bordo1993}。国际清算银行后来回顾说，那几年外国对美元商品的进口受到严格的数量限制，“美国持续的顺差显然让世界其他地方处境难以为继”\cite{bis36}。美国借给英国 37.5 亿美元，加拿大另借 12.5 亿美元，英国则批准布雷顿森林协定，并答应在一年以内恢复经常账户的可兑换\cite{bordo1993}。1945 年 12 月 6 日签订的美英贷款协定，1946 年 7 月 15 日获得批准\cite{bis36,bordo1993}。
% 据: bo_oeec, b68_turn, bo_loan, bo_loan_year, b48_dates, bo_loan_ratified

1947 年 7 月 15 日，英镑按约定恢复了经常账户交易的可兑换。挤兑在一个月内耗掉英国 10 亿美元的储备，8 月 20 日可兑换就中止了\cite{bordo1993}。照国际清算银行 1948 年年报的解释，根本的原因在于“经济结构与金融结构都还没有充分平衡，不足以恢复经常交易的完全可兑换”\cite{bis36}。美元的需求太大，伦敦市场面对的以英镑换美元的要求，规模是“战前从未见过的”\cite{bis36}。英国自己 1947 年的国际收支逆差还有近 7 亿英镑\cite{bis36}。英镑是英国的货币单位。
% 据: bo_1947, b48_structure, b48_dollarneeds, b48_deficit

基金组织在这些年里作用不大。成员国用本币向基金组织买外汇，叫提款（\secref{sec:warfinance-bw}）。1952 年以前虽有少量提款，国际收支的结构性援助大部分来自马歇尔计划与美国的其他援助（\chref{ch:reconstruction}）\cite{bordo1993}。法国 1948 年起实行多重汇率，基本进口按每美元约 214 法郎的官价，旅游与金融交易用浮动汇率，硬通货的实际汇率约 260 法郎\cite{bordo1993}。法郎是法国的货币单位。这样一来，在巴黎套算出来的美元与英镑的比价偏离了两者的平价，基金组织因此按协定第四条第 5 节谴责了法国，1952 年以前不许它动用基金的资源\cite{bordo1993}。法国有马歇尔援助，这些处罚效果不大\cite{bordo1993}。
% 据: bo_drawings, bo_france48, bt_fr_b, bo_censure, bo_denied, bo_noeffect

1949 年，危机又回到英镑。那年春天，一些原料与食品的价格大跌，“美国行情的下跌，带动了世界市场的下行”\cite{bis36}。英镑区里生产原料的国家对美国的出口减少，大大影响了英镑\cite{bis36}。英国的黄金与美元储备，从 4 月初到 9 月中减少了约 5.7 亿美元\cite{bis36}。照国际清算银行的说法，贬值也是为了不让英国的原料价格高于世界市场\cite{bis36}。
% 据: b50_slump, b50_exports, b50_reserves, b50_materials

贬值以前两个星期，财政大臣克里普斯（Stafford Cripps，1889—1952）还否认贬值在即\cite{bordo1993,loc36}。9 月 18 日，英镑贬值 30.5\%，降到 2.80 美元\cite{bordo1993}。原来的平价约合 4.03 美元（本书据 Bordo 表 1.2 的换算）\cite{bordo1993}。基金组织的工作人员早先劝过英国贬值，可这次贬的幅度比他们建议的大。按协定第四条第 5 节，贬值超过 10\% 要先与基金组织商量，给它 72 小时以上表示同意或反对。英国却只提前 24 小时通知了基金组织\cite{bordo1993}。几天以内，美元区与东欧以外的大部分货币也贬了值，“虽然幅度不尽相同”\cite{bis36}。随后降低平价的有 23 国\cite{bordo1993}。
% 据: bo_cripps, loc_cripps, bo_notice, b50_devalued, bt_uk

贬值的代价由谁承担，国际清算银行 1950 年的年报说得直白。在大多数情况下，贬值“等于用间接的办法降低实际工资”\cite{bis36}。比利时法郎 1949 年 9 月只贬了 12.3\%。按年报的说法，贬得多一些，也许能抵消工资成本的上升、缓和失业。可那样一来，“比利时的大部分工人至少要失去一部分加薪的好处”\cite{bis36}。
% 据: b50_belgium, b50_wages

瑞典与挪威走了另一条路。政府预算多发补贴，压住生活费用，又与工会等方面达成协议，“只要生活费用不超过一定的限度，就接受工资率不变”\cite{bis36}。补贴也有上限。1950 年 4 月，挪威的黄油涨价 23\%，人造黄油涨 125\%，面包涨了 40\% 以上\cite{bis36}。这是为了把生活费用补贴的总额压在 6 亿克朗以内，那是贬值以前的拨款\cite{bis36}。克朗是挪威的货币单位。英国的储备 1949 年 9 月底到 1950 年 3 月底增加了 5.6 亿美元，可年报也说，很难说其中有多少是贬值的功劳，美国的衰退并没有持续下去\cite{bis36}。在它看来，1949 年的危机虽然惊人，只是 1947 年以来靠马歇尔计划取得的长足进步中“一次暂时的挫折”\cite{bis36}。
% 据: b50_unions, b50_norway, b50_recovered, b50_unclear, b50_setback

这种钉住而又可以调整平价的安排，叫可调整的钉住（\chref{ch:depressionworld}）。它也给了投机者单向的赌注。平价只在根本性的失衡时才能改，货币当局总要拖到确信非改不可；到那时，投机者也确信了\cite{bordo1993}。货币贬了值，他们赢；不贬，他们只损失投机资金的利息\cite{bordo1993}。
% 据: bo_delay, bo_oneway

1949 年也是美元的转折点。国际清算银行 1968 年的年报写道，“回头看，1949 年外国货币的大批贬值，是美元的一个关键转折点”\cite{bis36}。在那以前，尽管美国的援助与官方贷款平均每年有 55 亿美元，美国的国际收支总是顺差，黄金储备上升\cite{bis36}。汇率调整以后，美国几乎年年逆差，黄金储备多数时候在下降\cite{bis36}。
% 据: b68_turn, b68_deficit

德国马克与日元也有了对美元的固定汇率（\chref{ch:reconstruction}）。德国马克是西德 1948 年起的货币单位，1949 年 9 月 19 日贬到每美元 4.20 马克\cite{bordo1993}。日元是日本的货币单位，1949 年 4 月 25 日定为每美元 360 日元\cite{bordo1993}。两国的官方平价都到 1953 年才正式定下\cite{bordo1993}。\tabref{tab:brettonwoods-parities}列出了几种主要货币的平价。
% 据: bt_de, bt_jp

\begin{table}[tbp]
\centering\small
\caption{几种主要货币的平价与变动，1946—1971 年}\label{tab:brettonwoods-parities}
\begin{tabularx}{\linewidth}{@{}l X@{}}
\toprule
国家 & 平价与变动 \\
\midrule
美国 & 1946 年起每盎司黄金 35 美元；1961 年 11 月成立伦敦黄金总库；1968 年 3 月 17 日起黄金分成两个市场；1971 年 8 月 15 日停止黄金兑换 \\
英国 & 1946 年每美元 0.248 英镑；1949 年 9 月 18 日贬到 0.357 英镑；1958 年 12 月 27 日事实上可兑换；1961 年 2 月 15 日按第八条可兑换；1967 年 11 月 18 日贬到 0.417 英镑 \\
法国 & 1946 年每美元 119.1 法郎；1948 年 1 月 26 日起平价中止；1949 年 9 月 20 日 350 法郎；1957 年 8 月 11 日 420 法郎；1958 年 12 月 29 日 493.7 法郎；1960 年起新法郎 1 元合旧法郎 100 元；1969 年 8 月 10 日贬到 5.55 新法郎 \\
西德 & 1948 年 6 月 20 日发行德国马克，每美元 3.33 马克；1949 年 9 月 19 日贬到 4.20 马克；1961 年 3 月 6 日升到 4 马克；1969 年 10 月 26 日升到 3.66 马克；1971 年 5 月 9 日浮动 \\
日本 & 1949 年 4 月 25 日定为每美元 360 日元；1953 年 5 月 11 日定为官方平价；1964 年 4 月 1 日按第八条可兑换 \\
\bottomrule
\end{tabularx}
\datasource{Bordo（1993）表 1.2，据页面图像读出；原表的来源是基金组织的各种出版物。}
\end{table}
% 据: bt_us, bt_uk, bt_fr, bt_de, bt_jp, bt_source

\section{回到可兑换}\label{sec:brettonwoods-convertibility}

战后头几年，照国际清算银行的说法，国际贸易与支付大都被困在“双边主义的紧身衣”里\cite{bis36}。各国尽量与每个贸易伙伴分别平衡收支，免得动用国际储备。在欧洲，“一国赚到的另一国货币，一般只能在那个国家花”\cite{bis36}。1947 年 11 月，几个欧洲国家签订了多边抵消的协定，这是走向更自由的支付的第一步\cite{bis36}。1947 年到 1950 年初，西欧各国之间的贸易量增加了近一半，西欧对世界其他地方的出口增加了约三分之一\cite{bis36}。
% 据: b59_straitjacket, b59_nov47, b50_trade

从 1950 年 7 月起，欧洲支付同盟提供了一个“每月完全结清的简单机制”，让西欧各国之间的收支可以多边抵消\cite{bis36}。每月月底，每个成员对其他所有成员的净借方或净贷方头寸，一律与同盟清算，国际清算银行作代理人\cite{bordo1993}。记账单位是美元，结算用美元、黄金或信贷\cite{bordo1993}。美国提供了 3.5 亿美元的初始周转资金\cite{bordo1993}。一国对一个伙伴的逆差，可以用对另一个伙伴的顺差抵消，用不着为每一笔双边逆差去找美元（本书的解读）。
% 据: b59_epu, bo_epu_month, bo_epu350

1958 年 12 月底，同盟的成员除希腊、冰岛与土耳其以外，连同芬兰，同时对非居民恢复了经常账户的可兑换\cite{bis36}。也就是说，外国人在经常交易里赚到的这些货币，可以自由换成美元。同盟的协定在 12 月 27 日营业结束时终止\cite{bis36}。照国际清算银行的说法，这取消了欧洲内部结算中自动的官方信贷，也意味着“世界主要外汇市场的统一”\cite{bis36}。
% 据: b59_conv, b59_end, b59_winding

1961 年 2 月 15 日，比利时、法国、西德、爱尔兰、意大利、卢森堡、荷兰、秘鲁、瑞典与英国 10 国接受了协定第八条的义务\cite{bis36,bordo1993}。第八条要求它们“避免限制经常的对外支付，避免多重汇率与歧视性的货币做法”\cite{bis36}。日本到 1964 年 4 月 1 日才接受\cite{bordo1993}。照 Garber 的说法，体系在完全可兑换下真正按设想运转的“全盛期”，只有 1959—1968 年\cite{garber1993}。
% 据: b61_art8, b61_art8b, b61_art8c, bo_art8, bt_jp, ga_heyday

基金组织的资源也在增加。1960 年各国的份额增加 50\%，1966 年又增加 25\%\cite{bordo1993}。国际清算银行 1961 年的年报记载，基金组织第一次有将近一半的新提款用的是美元以外的货币，共 1.31 亿美元\cite{bis36}。接受第八条的国家，提款可以用它们的货币偿还，受第八条约束的成员已有 21 个\cite{bis36}。
% 据: bo_quotas, b61_drawings, b61_repaid

可兑换以后，钱有了更多的出入口，利率的差别开始起作用。西德 1959 年取消了资本管制，1960 年又重新限制资本流入\cite{bordo1993}。1960 年，西德在提高短期利率，美国却在降低\cite{bis36}。西德的企业转到利率低的外国去借钱，外国资本为了高收益流进来，出口的货款提前收，进口的货款推迟付\cite{bis36}。照国际清算银行的说法，有了可兑换，货币政策这张网上的“洞太大了”\cite{bis36}。西德限制非居民在西德的短期投资，收效甚微\cite{bis36}。
% 据: bo_controls_de, b61_1960, b61_rates, b61_shift, b61_influx, b61_holes, b61_limited

1960 年 11 月与 1961 年 1 月，需求并没有放缓，西德却两次降低贴现率\cite{bis36}。照国际清算银行的说法，这“实际上承认了”在当时的情况下用不了更高的利率，从这里到决定马克升值，“步子并不大”\cite{bis36}。
% 据: b61_cuts, b61_admission

国际清算银行列出了西德当局可以采取的三类办法，彼此并不排斥\cite{bis36}。一是放松信贷，让需求上升，用更多的进口与私人资本输出吸收顺差，成本与物价也许随之上升。二是让马克升值，以便放松国内信贷而不推高成本与物价。三是由政府对外援助与贷款，抵消出口的顺差\cite{bis36}。荷兰的情况相似，同时达到内部与外部的均衡“越来越难”\cite{bis36}。
% 据: b61_three, b61_inner

1961 年 3 月 6 日与 7 日，马克与荷兰盾先后升值 5\%，马克升到每美元 4 马克\cite{bis36}。荷兰盾是荷兰的货币单位。两国官方的声明都强调这是一次性的调整。可是升值的幅度不大，此前四个月黄金对美元又一直有可观的溢价，照国际清算银行的说法，市场的第一反应“毫不奇怪地反映了颇为不同的看法”\cite{bis36}。3 月的第二周，西德的货币储备又增加了 2 亿多美元，瑞士增加了近 2.5 亿美元，英国的储备却急剧减少\cite{bis36}。
% 据: b61_reval, b61_dm4, b61_onceforall, b61_inflows

照 Bordo 的概括，西德在整个布雷顿森林时期都抗拒调整，它的货币当局认为，国际货币体系的关键问题是“从国外输入的通胀”\cite{bordo1993}。资本一旦流动起来，固定的汇率与自己的货币政策就难以兼得，这个道理留给\chref{ch:mundell}。
% 据: bo_de_resist

从数字看，可兑换以后的十年运转得最好。照 Bordo 的比较，1946—1958 年可兑换以前的时期，名义与实际变量都远不如后来稳定\cite{bordo1993}。可这些数字“没有告诉我们”它为什么运转得这么好，是环境有利，还是没有遇到加重的冲击\cite{bordo1993}。
% 据: bo_preconv, bo_whywell

\section{特里芬难题}\label{sec:brettonwoods-triffin}

1950 年代的美国逆差，起初是受欢迎的。国际清算银行 1968 年回顾说，这些逆差“大体上被看作建设性的，确实大大帮助了欧洲在迅速重建的同时积累储备”\cite{bis36}。1950—1956 年，按官方储备交易计算，美国的逆差平均每年约 12 亿美元\cite{bis36}。对外援助从 1949 年的 56 亿美元减到 1956 年的 24 亿美元，美国在海外的防务支出却大大增加\cite{bis36}。
% 据: b68_constructive, b68_5056, b68_aid, b68_defence50s

1957 年中东打仗，苏伊士运河关闭，欧洲国家在西半球大量买油\cite{bis36}。美国这一年的贸易顺差有 60 亿美元，国际收支顺差 11 亿美元，年报说那是“暴风雨前的一阵阳光”\cite{bis36}。1958 年，美国的国际收支转为 29 亿美元的逆差\cite{bis36}。
% 据: b68_1957y, b68_suez, b68_1957

1960 年，特里芬（Robert Triffin，1911—1993）出版了《黄金与美元危机》\cite{bordo1993,gnd36}。照 Bordo 的概括，特里芬论证流动性短缺，主要的根据是黄金储备不足\cite{bordo1993}。储备增长的缺口，英镑余额补不上，它已在减少。基金组织只能补上一部分，其余要靠美国的国际收支逆差\cite{bordo1993}。逆差持续下去，美国的货币黄金无论绝对地还是相对于未偿的美元负债都会减少，最后出现兑换的危机\cite{bordo1993}。美国若在那以前消除逆差，又会造成国际流动性的严重短缺，世界面临通缩\cite{bordo1993}。
% 据: bo_ref_triffin, gnd_triffin, gnd_triffin_born, gnd_triffin_died, bo_tr_gold, bo_tr_rest, bo_tr_crisis

按 Bordo 引的原话，美元余额“不能指望”长期大量地解决世界的流动性问题\cite{bordo1993}。特里芬自己的办法，是回到凯恩斯（John Maynard Keynes，1883—1946）的国际清算同盟方案\cite{bordo1993,loc36}。现有的储备都换成国际货币，基金组织充当世界的中央银行\cite{bordo1993}。这本账可以写成三个式子。
% 据: bo_triffin, bo_tr_solution, loc_keynes

\begin{keyeq}
R_t=G_t+D_t,\qquad \Delta D_t=gR_{t-1}-\Delta G_t,\qquad c_t=\frac{G^{\mathrm{US}}_t}{D_t}.\label{eq:brettonwoods-triffin}
\end{keyeq}
$R$是各国货币当局的储备，$G$是其中的黄金，$D$是其中的美元。它与\chref{ch:interwar}的式~\eqref{eq:interwar-pyramid}里的$R=G+F$是同一本账，只是外汇换成了美元。储备的需要每年增长$g$，黄金只增加$\Delta G$，差额由美元补上，$\Delta D$就是美国对外国货币当局的逆差。$c$是美国的黄金$G^{\mathrm{US}}$对外国持有的美元的覆盖率。$c$大于 1，所有的美元都换得到黄金；$c$小于 1，大家一起来换，黄金就不够了。$g$越大，$\Delta G$越小，$c$跌破 1 就越早（本书的模型）。

\begin{example}{储备的增长与黄金的覆盖率}\label{ex:brettonwoods-triffin}
设第 0 年各国的储备$R_0=40$，其中黄金$G_0=30$，美元$D_0=10$；美国的黄金$G^{\mathrm{US}}=20$不变。储备的需要每年增长 5\%，黄金每年只增加 0.5，数字为教学输入。（1）第 1 年美元要增加多少？（2）写出$D_t$。（3）哪一年$c$跌破 1？（4）若黄金每年增加 1.5 呢？
\solution
（1）$R_1=40\times1.05=42$，$\Delta D_1=0.05\times40-0.5=1.5$，$D_1=11.5$，$c_1=20/11.5\approx1.74$。（2）$R_t=40\times1.05^t$，$G_t=30+0.5t$，所以$D_t=40\times1.05^t-30-0.5t$。（3）第 5 年$D_5\approx51.05-32.5=18.55$，$c_5\approx1.08$；第 6 年$D_6\approx53.60-33=20.60$，$c_6\approx0.97$，第 6 年跌破 1。（4）$G_t=30+1.5t$。第 9 年$D_9\approx62.05-43.5=18.55$，$c_9\approx1.08$；第 10 年$D_{10}\approx65.16-45=20.16$，$c_{10}\approx0.99$，推迟到第 10 年。黄金增加得快一些，只能推迟，避免不了：储备按比例增长，黄金只按固定的量增加，美元的份额迟早越来越大。这里还假定美国的黄金不变。外国货币当局拿美元换黄金，$G^{\mathrm{US}}$与$D$同时减少：$c$大于 1 时，兑换反而让覆盖率上升；$c$小于 1 时，兑换让它跌得更快。
\end{example}

数字正是这样走的。国际清算银行 1968 年说，“黄金存量上升是实力的标志，下降至少应该引起疑问”\cite{bis36}。1950—1956 年，美国的黄金平均每年减少约 4 亿美元\cite{bis36}。照国际清算银行的说法，原因并不在人们怀疑美元将来能否兑换黄金。新增的货币黄金平均每年 4.4 亿美元，只够外国货币当局净需要的一半多一点，其余的它们向美国财政部取得\cite{bis36}。可年报也说，“与官方负债的上升合起来看，黄金的流失不是好兆头”\cite{bis36}。到 1965 年，新黄金的供应创下近 20 亿美元的纪录，私人买走了约 16 亿美元，留给官方的只有 4 亿美元\cite{bis36}。
% 据: b68_indicator, b68_newgold, b68_distribution, b66_record, b66_private

\begin{example}{美国的数字}\label{ex:brettonwoods-us}
按美国《统计摘要》各版的表，年底美国的黄金存量、对外国的短期负债与其中对外国官方机构的负债如下（亿美元）：1955 年 217.53、117.20、69.53；1960 年 178.04、172.60、102.12；1967 年 120.65、292.32、140.27\cite{statab36}。数据来自联邦储备委员会与《联邦储备公报》\cite{statab36}。（1）算出三年的两种覆盖率。（2）对外国的负债与对官方的负债，各在哪一年超过黄金？（3）1955—1967 年，黄金减少了多少，对外国的负债增加了多少？
\solution
（1）1955 年$217.53/117.20\approx1.86$，$217.53/69.53\approx3.13$；1960 年$178.04/172.60\approx1.03$，$178.04/102.12\approx1.74$；1967 年$120.65/292.32\approx0.41$，$120.65/140.27\approx0.86$。（2）1961 年底对外国的负债是 187.81 亿美元，黄金 169.47 亿美元，覆盖率约 0.90，第一次低于 1\cite{statab36}。1966 年底对官方的负债是 125.39 亿美元，黄金 132.35 亿美元，覆盖率约 1.06，1967 年降到约 0.86\cite{statab36}。（3）黄金减少了约 96.9 亿美元，对外国的负债增加了约 175.1 亿美元。美元增加的差不多是黄金减少的两倍，式~\eqref{eq:brettonwoods-triffin}里$D$上升、$G^{\mathrm{US}}$下降，两头一起压低$c$（本书的计算）。
\end{example}
% 据: sa72_head, sa72_gold, sa72_foreign, sa72_official, sa68_source, sa64_gold, sa64_foreign, sa64_fn2, sa64_official, sa68_gold, sa68_foreign, sa68_official

\figref{fig:brettonwoods_gold}画出了 1955—1971 年的数字。照 Bordo 的叙述，到 1959 年，美国的货币黄金已与对外美元负债的总额相等，其余各国的货币黄金也超过了美国\cite{bordo1993}。按《统计摘要》1964 年版，1959 年底包括国际组织在内的对外短期负债约 194 亿美元，黄金约 195 亿美元，两者确实差不多\cite{statab36}。官方的负债什么时候超过黄金，说法不一。照 Bordo 的说法是 1964 年，照 Garber 的说法是 1960 年代后半期，《统计摘要》的数字是 1967 年底\cite{bordo1993,garber1993,statab36}。国际清算银行 1968 年的年报算过整段时期的账\cite{bis36}。从 1950 年初到 1968 年 4 月底，美国的黄金储备从 246 亿美元降到 107 亿美元，后者还包括欠基金组织的 10 亿美元\cite{bis36}。同一时期，对外国官方机构的负债从不到 40 亿美元增加到 150 亿美元以上\cite{bis36}。
% 据: bo_1959, sa64_total, sa64_gold, bo_1964, ga_wrong, sa64_official, sa68_official, sa72_official, b68_gold

统计本身也在变。《统计摘要》对官方的负债，只包括外国中央银行、政府及其机构、国际清算银行与欧洲基金的美元\cite{statab36}。照 Garber 的说法，美国国内居民的债权也远远超过了美国的黄金。美国公民拥有黄金虽然非法，这些债权却可以转到国外，很快拿来挤兑，可当时的讨论没有理会它们\cite{garber1993}。国际清算银行 1970 年指出，美国报告的对外国官方的负债，只包括存在美国境内银行的美元\cite{bis36}。照它的推断，1969 年“存在美国以外银行的美元储备必定大大增加了”，所以图上那一年的官方负债减少了\cite{bis36}。
% 据: sa64_fn4, sa72_fn, ga_citizens, b70_eurodollar, b70_outside

\fig{brettonwoods_gold}{美国的黄金与对外国的美元负债，1955—1971 年（年底，亿美元）。(a) 黄金存量、对外国的短期负债及其中对外国官方机构的负债。(b) 黄金对两种负债的覆盖率，虚线是 1。1956—1958 年缺数。对外国的负债 1961 年、对官方的负债 1967 年超过了黄金；1969 年对官方的负债减少，照国际清算银行的推断，是因为这项统计只包括外国官方存在美国境内银行的美元。数据：美国《统计摘要》1964、1968、1972 年版，表 1211、1204、1271（联邦储备委员会，《联邦储备公报》）；国际清算银行 1970 年年报。}{fig:brettonwoods_gold}

\begin{compare}{布雷顿森林体系与古典金本位}{古典金本位}{布雷顿森林体系}
\cmprow{什么按固定价格换黄金}{英镑当时可以按固定的价格兑换黄金（\chref{ch:goldstandard}）}{只有美元，每盎司 35 美元，只对外国的货币当局\cite{bordo1993}}
\cmprow{各国的储备}{黄金；印度等地把储备存在伦敦（\chref{ch:periphery}）}{黄金与美元，美元越多，覆盖率$c$越低（式~\eqref{eq:brettonwoods-triffin}）}
\cmprow{平价}{停兑以后按旧的平价恢复，英国 1821 年就是这样（\secref{sec:bullion-resumption}）}{可以调整，英国、法国、西德都改过（\tabref{tab:brettonwoods-parities}）}
\cmprow{谁先调整}{伦敦用高利率向世界要黄金（\chref{ch:goldstandard}）}{储备货币国美国不必为国际收支调整国内经济\cite{bordo1993}}
\cmpsum{金本位下，黄金是最后的储备。布雷顿森林体系把美元放在黄金与各国之间，美元既是别国的储备，又是美国要拿黄金兑付的负债。}
\end{compare}
% 据: bo_pegged, bt_us, bo_asym

\begin{pitfall}[布雷顿森林体系下人人都能拿美元换黄金]
能按 35 美元换黄金的，只有外国的货币当局\cite{bordo1993}。美国居民拥有黄金是非法的\cite{garber1993}。1960 年秋天金价上涨以后，美国又规定本国公民在国外购买或持有黄金也是非法\cite{bis36}。私人要买黄金，只能去伦敦这样的市场，那里的价格可以离开 35 美元，1960 年 10 月就涨到过 40 美元\cite{bordo1993}。黄金总库卖出黄金，就是为了把这个市场的价格钉在 35 美元\cite{bordo1993}。美元的可兑换，是对各国货币当局的承诺；1968 年 3 月以后，官方持有的黄金只在货币当局之间转移（\secref{sec:brettonwoods-defence}）。
\end{pitfall}
% 据: bo_pegged, ga_citizens, b62_measures, bo_oct60, bo_pool_nov, fr191_communique

\section{保卫美元}\label{sec:brettonwoods-defence}

1960 年秋天，伦敦的金价离开了美国的官价。美国财政部卖出黄金的价格约为每盎司 35.09 美元。9 月中到 10 月中，伦敦的报价却在 35.20 到 35.25 美元之间，比从纽约买进黄金、运到伦敦转卖的保本价还高几美分\cite{bis36}。10 月，伦敦的自由价格涨到了 40 美元\cite{bordo1993}。10 月 31 日，总统候选人肯尼迪（John F. Kennedy，1917—1963）在演说里承诺，“如果当选总统，我不会让美元从现在的汇率上贬值”\cite{bordo1993,loc36}。国际清算银行 1962 年说，这次风波表明，不受遏制的黄金投机“可能威胁国际货币体系”\cite{bis36}。美国随后规定本国公民在国外购买或持有黄金为非法；每月在国际清算银行开会的各国中央银行也非正式地约定，市场有压力时不去买黄金\cite{bis36}。
% 据: b61_selling, b61_london, bo_oct60b, bo_oct60, bo_kennedy, loc_kennedy, b62_threat, b62_measures

伦敦黄金总库 1961 年 11 月正式成立。美国与另外 7 家中央银行约定买卖黄金，把价格钉在每盎司 35 美元\cite{bordo1993}。按约定，美国提供 50\% 的黄金\cite{garber1993}；按 Bordo 引的数字，实际卖出的黄金有 60\% 来自美国\cite{bordo1993}。另外 7 家虽然提供了稳定金价所需黄金的 40\%，却在总库以外把手里的美元余额拿到美国财政部换成黄金，补回自己的货币黄金\cite{bordo1993}。黄金从美国流出的路，并没有因总库而堵上。照国际清算银行的说法，1961—1965 年美国黄金减少得较少，部分是因为别国在当时紧张的情况下不愿把美元储备换成黄金\cite{bis36}。
% 据: bo_pool_nov, bt_us, ga_pool50, bo_pool60, bo_pool40, b62_pool, b68_1961_65, b68_reluctant

总库起初管用。1965 年私人买金最多的时候，总库暂时投入市场的黄金约 2.5 亿美元，到年底又全部买了回来\cite{bis36}。国际清算银行由此得出结论，即使公众有很强的投机冲动，总库也“能够轻易地控制市场”\cite{bis36}。两年以后，情形就不同了。
% 据: b66_pool, b66_recouped, b66_control

其他的办法一样样加了上来。1961 年起，美国与其他中央银行签订互换协定，双方互相提供双边的信贷额度\cite{bordo1993}。为了偿还互换，财政部发行以外币计值的债券，叫罗萨债券，卖给外国的官方机构\cite{bordo1993,bis36}。1963 年年中起，美国对本国居民购买外国证券征收利息平衡税。1965 年又推出自愿计划，限制银行对外贷款与企业的对外直接投资\cite{bis36}。利息平衡税 1964 年 9 月 2 日才正式立法，1967 年 7 月 31 日又延长到 1969 年 7 月 31 日\cite{frus36}。
% 据: bo_swaps, bo_roosa, b68_roosa, b68_iet, fr135_first, fr135_iet

1961 年 12 月，十个工业国议定，基金组织可以向它们借入 60 亿美元，“去帮助本集团里货币受到短期资本外逃威胁的成员”，这就是借款总安排\cite{bis36}。这十国就是后来的十国集团。安排 1962 年 10 月开始运作\cite{garber1993}。照 Garber 的说法，借款总安排、互换与后来由基金组织分配的特别提款权，都是在原有信贷额度之上“临时搭起来的”新的信贷额度\cite{garber1993}。
% 据: b62_gab, ga_t91, ga_gab

这些办法搭起了框架，1968 年 1 月 1 日美国又宣布了更广、更严的管制\cite{bis36}。美国公司的对外直接投资，从自愿的准则改为强制的管制\cite{bis36}。国际清算银行说，这是“一条危险的路，可能让世界离开过去十五年左右达到的自由的贸易与支付做法”\cite{bis36}。管制的理由是越南战争，以及战时难以恢复真正的均衡。年报接着问，“和平回来以后，管制就能拆掉吗？”\cite{bis36}
% 据: b68_jan1968, b68_programme, b68_dangerous, b68_peace

国内的政策却很少动。照国际清算银行 1968 年的概括，美国给自己定了三条界限\cite{bis36}。一是大幅削减驻外的防务力量，“在当时的国际政治形势下被认为行不通”。二是不宜改变美元的金价。三是低的长期利率被认为是加快国内经济增长所必需的\cite{bis36}。货币政策只用了“扭转操作”，托住短期利率而不影响长期利率\cite{bis36}。
% 据: b68_limits, b68_twist

年报指出，1966 年为抑制国内通胀实行紧的货币政策，美国资本账户的逆差就大大减少，可见用这种办法“本可以更早取得明显的效果”\cite{bis36}。驻外的军队、金价与国内的增长都不动，账就落在资本管制与外国的中央银行身上（本书的解读）。1965 年起，美国的货币增长加快，部分是为了给越南战争筹资\cite{bordo1993}。
% 据: b68_1966, bo_vietnam

英镑先撑不住。照 Bordo 的概括，英国的增长比贸易伙伴慢，通胀却更高，稳定政策又走走停停，英镑的竞争地位受到威胁；1957 年就有过一次大的危机与基金组织的救援\cite{bordo1993}。1964 年 10 月 16 日工党当选时，储备迅速减少。新政府拒绝贬值，10 月 26 日宣布对进口加征附加税\cite{bordo1993}。11 月 25 日，十国集团与基金组织为英镑安排了 40 亿美元的救援\cite{bordo1993}。1966 年 7 月 20 日，英国实行大规模的紧缩，1967 年初产出下降、失业上升，政策才又放松\cite{bordo1993}。六日战争中苏伊士运河关闭，10 月又有码头罢工，11 月英镑遭到投机冲击；30 亿美元的救援没能挡住潮流，11 月 18 日英镑贬值 14.3\%，降到 2.40 美元\cite{bordo1993}。这次贬值实际上结束了英镑作为储备货币的角色\cite{bordo1993}。
% 据: bo_stopgo, bo_uk_problem, bo_uk1964, bo_rescue64, bo_1966, bo_uk1966, bo_shocks, bo_1967, bo_endreserve

11 月 20 日，财政大臣卡拉汉（James Callaghan，1912—2005）在下院说明了贬值与配套的措施\cite{hansard1967,loc36}。贬值是他向内阁提议、内阁接受的\cite{hansard1967}。银行利率提高到 8\%，银行贷款除了给出口商等优先借款人以外都受到限制\cite{hansard1967}。他说，“我们从国外买的东西要多花钱，我们卖到国外的东西对外国人会便宜”，未来几个月物价会有小而明显的上涨\cite{hansard1967}。政府会在适当的时候采取措施，“保护社会上最脆弱的人群免受困苦”\cite{hansard1967}。
% 据: h67_speaker, loc_callaghan, h67_cabinet, h67_rate, h67_measures, h67_prices, h67_rise, h67_vulnerable

要求首先摆给了工资。贬值的主要目的，是让国际收支得到持久而可观的改善，每年至少 5 亿英镑；物价不可避免的上涨若被普遍的加薪抵消，这一切就会白费\cite{hansard1967}。他说这正是他 7 月的意思，并把当时在下院的话重念了一遍：主张贬值的人，是在要求降低“这个国家每一个工人阶级成员的工资水平与实际工资”\cite{hansard1967}。若交货误期，质量不够，大的加薪要求提了出来又被答应，竞争上的好处就会“白白浪费”\cite{hansard1967}。生产率真正大幅提高的，可以在工资谈判里体现；股息也会受到严格的监视\cite{hansard1967}。
% 据: h67_500, h67_wages, h67_meant, h67_july, h67_claims, h67_productivity, h67_dividends

照他的估计，不贬值的话，1968 年的个人消费会增加 3\%，国际收支却会有“很大的、不可接受的”逆差\cite{hansard1967}。贬值以后，本来要增加的个人消费大部分转到出口上去，他说，“这是我们必须付出的代价”\cite{hansard1967}。公共支出的预期增长削减 4 亿英镑，明年预算里的防务削减 1 亿英镑\cite{hansard1967}。持有英镑的英镑区国家，手里英镑的外汇价值也减少了\cite{hansard1967}。
% 据: h67_1968, h67_consumption, h67_price, h67_cuts, h67_defence, h67_holders

他也说，各国提供的信贷“不是给我们的贷款，不是让我们靠它过日子的”\cite{hansard1967}。贬值是为了把生产转向出口，“我们要的是出口带动的增长”\cite{hansard1967}。出口一增加，失业就会开始减少，他预计冬天结束以后失业会稳步下降\cite{hansard1967}。
% 据: h67_loan, h67_exportled, h67_winter

结果并不如意。贬值以后，进口的数量急剧上升，出口虽然大幅增加，国际收支却改善不多\cite{bis36}。1968 年 11 月与 1969 年 4 月，紧缩又两次收紧，加税，加强分期付款的限制，压低银行贷款的上限\cite{bis36}。英镑区国家此后越来越多地改持美元。1968 年在巴塞尔达成协议，每个成员以一定比例的储备持有英镑，换来对超过总储备 10\% 的英镑的美元担保\cite{bordo1993}。
% 据: b70_uk, b70_uk_tight, bo_basle68

\begin{example}{英镑贬值的账}\label{ex:brettonwoods-sterling}
1967 年 11 月，英镑从 2.80 美元降到 2.40 美元\cite{hansard1967}。设英国一年出口按英镑计价 100 英镑，进口按美元计价 280 美元，贬值以前也合 100 英镑；各自的本币价格不变，数字为教学输入。（1）贬值以后，一美元合多少英镑？英镑的美元价值降了百分之几？（2）出口与进口的数量都不变，以英镑计的贸易差额是多少？（3）若出口的数量增加 15\%，进口的数量减少 10\%，差额又是多少？用式~\eqref{eq:external-bop}说明它对储备的意义。
\solution
（1）一美元从$1/2.80\approx0.357$英镑变成$1/2.40\approx0.417$英镑，贵了约 16.7\%；英镑的美元价值降了$(2.80-2.40)/2.80\approx14.3\%$，正是\tabref{tab:brettonwoods-parities}里的两个数。（2）出口仍收 100 英镑，在外国只卖$100\times2.40=240$美元；进口要付$280/2.40\approx116.7$英镑，差额从 0 变成约$-16.7$英镑。卡拉汉说的“多花钱”，先落在进口上。（3）出口$100\times1.15=115$英镑，进口$116.7\times0.9\approx105$英镑，顺差约 10 英镑。按式~\eqref{eq:external-bop}，资本账户与误差两项不变，经常账户改善多少，金融账户的净增加就多多少；私人的对外资产与负债也不变的话，多出来的全是官方储备。改善要靠数量的反应；国际清算银行看到的却是，1967 年 11 月贬值以后进口的数量急剧上升，出口虽然大增，国际收支改善不多\cite{bis36}。
\end{example}
% 据: h67_rate, h67_prices, bt_uk, bo_1967, b70_uk

英镑贬值以后，压力转到了黄金上。照 Bordo 引的 Johnson（1968）的看法，英国若在 1964 年而不是 1967 年贬值，1968 年对黄金与美元的投机冲击也许不会发生\cite{bordo1993}。市场上的净销售通常每天三四吨，贬值以前几周常常翻倍\cite{bis36}。人们这时才知道，法国从 6 月起已不再参加总库对市场供应黄金\cite{bis36}。11 月 26 日，总库的活跃成员在法兰克福开会\cite{bis36}。这一年第四季度，西方的货币黄金减少了 14 亿美元\cite{bis36}。1967 年 12 月到 1968 年 3 月，总库卖掉了 30 亿美元的黄金，其中美国 22 亿美元\cite{bordo1993}，约占 73\%（本书的计算）。3 月 14 日一天就卖出 4 亿美元\cite{garber1993}。
% 据: bo_fn52, b68_tons, b68_france_pool, b68_frankfurt, b68_q4, bo_pool_lost, ga_march

3 月 15 日星期五，伦敦的黄金市场关闭，巴黎市场的金价超过了 44 美元\cite{bis36}。3 月 16 日与 17 日，总库 7 国的中央银行行长在华盛顿开会\cite{frus36}。公报说，官方持有的黄金“今后只应用于货币当局之间的转移”，他们决定不再向伦敦或任何别的黄金市场供应黄金\cite{frus36}。他们也“不再把黄金卖给货币当局，去补它们在私人市场上卖掉的黄金”，总库以外的那条路也堵上了\cite{frus36}。既然特别提款权就要设立，他们认为现有的货币黄金已经够用，也不再从市场买黄金\cite{frus36}。
% 据: b68_paris, fr191_meeting, fr191_communique, fr191_replace, fr191_buy, fr191_35, bo_disbanded

公报还写下了各方的承诺。英国政府决心“做一切必要的事”来结束国际收支逆差，大多数欧洲国家的政府打算推行“鼓励国内扩张、又与经济稳定相容的”货币与财政政策\cite{frus36}。立即可以给英国的信贷，包括基金组织的备用安排，增加到 40 亿美元\cite{frus36}。备用安排是事先议定、需要时可以随时提用的贷款。
% 据: fr191_ukall, fr191_uk

黄金从此有了两个市场。官方交易的价格固定在每盎司 35 美元，私人交易的价格“由供求自由决定”\cite{frus36}。伦敦市场关闭了两周，4 月 1 日重开时金价是 38 美元，5 月下半月超过了 42 美元\cite{bis36}。1968 年 11 月到 1969 年 3 月，金价又从约 39 美元涨到将近 44 美元\cite{bis36}。南非卖出黄金以后，到 1969 年底，私人市场的金价才回到 35 美元左右\cite{garber1993}。
% 据: fr191_two, b68_closed, b68_42, b70_gold39, ga_42

新的储备也在准备之中。1967 年 9 月 29 日，基金组织在里约热内卢的年会批准了特别提款权方案的纲要\cite{frus36}。特别提款权由基金组织按份额记给各成员，没有黄金作后盾\cite{bordo1993}。一个特别提款权定义为等于一个黄金美元\cite{bordo1993}，按每盎司 35 美元的金价，就是三十五分之一盎司黄金（本书的换算）。1968 年 3 月 29 日与 30 日，十国集团在斯德哥尔摩开会\cite{frus36}。法国财政部长德勃雷（Michel Debré，1912—1996）在会上表示，修正案的草案远远超出了纲要，会让特别提款权更像储备账户，而不像补充的信贷\cite{frus36,loc36}。法国代表团保留了立场\cite{frus36}。
% 据: fr193_rio, bo_sdr_rio, bo_sdr_quota, bo_sdr_def, fr193_date, fr193_debre, loc_debre, fr193_reserved

照 Bordo 的说法，作为反对美国霸权的斗争的一部分，法国实际上迫使方案加上一条：只有美国的国际收支逆差消除以后，才能启动\cite{bordo1993}。修正案 1969 年 7 月生效\cite{garber1993}。美国 1968 年与 1969 年都是国际收支顺差，方案在 1970 年 1 月 1 日启动，三年分配 95 亿美元\cite{bordo1993}。
% 据: bo_sdr_france, ga_t91, bo_sdr_1970, bo_sdr_95

\begin{case}{法国 · 1958—1968 · 把美元换成黄金}
法国自己的平价也改过几次。1957 年 8 月，它对原来的汇率加上 20\% 的溢价或附加费，事实上贬了值，1958 年 6 月 21 日正式化\cite{bis36}。1958 年底又贬值将近 15\%，同时放开 90\% 的进口\cite{bis36}。新的美元平价是 493.706 法郎，12 月 29 日正式向基金组织申报；在那以前，法国自 1948 年 1 月 26 日起在基金组织一直没有申报的平价\cite{bis36}。1960 年起，1 新法郎合 100 旧法郎\cite{bordo1993}。
% 据: b59_fr57, b59_fr_end58, b59_fr58, b59_fr_parity, b59_fr_imf, bt_fr

这一次的代价，国际清算银行 1959 年的年报也记了下来。1958 年底的措施主要靠取消大量补贴来压预算赤字\cite{bis36}。有些食品的补贴全部取消，几种农产品的补贴部分取消，国有企业的收费也提高了\cite{bis36}。所有把工资或农产品价格与物价指数挂钩的协议都取消了，只有最低工资仍与“179 种商品指数”挂钩\cite{bis36}。照年报的说法，1957 年与 1958 年为理顺价格结构采取的各项措施，本来就会推高物价；1958 年，食品价格因补贴减少而急剧上涨，法国的生活费用涨了 8.9\%\cite{bis36}。
% 据: b59_fr_cuts, b59_fr_food, b59_fr_index, b59_fr_expected, b59_fr_prices

法国 1962 年第一季度开始买黄金时，美国的黄金存量是 168 亿美元；到 1966 年第三季度减少的 33 亿美元里，28 亿美元是法国直接买走的\cite{garber1993}。法国持续的购买到 1967 年 10 月才停止\cite{garber1993}。照 Bordo 的叙述，1965 年起法国有计划地把美元负债换成黄金\cite{bordo1993}。
% 据: ga_france, ga_france2, ga_t91, bo_convert

照 Gourinchas 等人的说法，财政部长吉斯卡尔·德斯坦（Valéry Giscard d'Estaing，1926—2020）1965 年创造了“过度特权”这个说法\cite{gourinchas2010,loc36}。照 Lelart 的叙述，他 1964 年 9 月在东京的基金组织年会上就用过它\cite{lelart2018}。大约同时，经济顾问吕夫（Jacques Rueff，1896—1978）把美元叫作“回旋镖货币”\cite{gourinchas2010,gnd36}。照 Gourinchas 等人的解释，美国的大笔逆差没有带来相应的黄金流失，因为债权国又把赚到的美元投回了美国\cite{gourinchas2010}。法国开出的药方是把金价翻一番，这正是 1934 年以来黄金实际价格下降的幅度\cite{bordo1993}。照 Bordo 对吕夫（1967）的概括，各国货币黄金重估的收益足以清偿未偿的美元与英镑余额，美国恢复国际收支平衡以后，世界就可以回到完全运转的古典金本位\cite{bordo1993}。
% 据: im_giscard, loc_giscard, le_tokyo, im_rueff, gnd_rueff, gnd_rueff_born, gnd_rueff_died, bo_convert, bo_rueff

黄金分成两个市场以后不久，法国出了政治问题，美国借此在此后一年里买进了将近 10 亿美元的法国黄金\cite{garber1993}。
\end{case}
% 据: ga_buyback

\section{三个争论}\label{sec:brettonwoods-debate}

体系为什么没有撑下去，至今没有共识。照 Eichengreen 的概括，在美国国家经济研究局的会议上，至少提出了六种解释\cite{eichengreen1993}。它们是美国与外国货币政策的差别，财政政策的差别，逆差国不贬值，顺差国不升值，美国国际竞争地位的长期下降，以及体系结构上的缺陷，特里芬的流动性两难就是一例\cite{eichengreen1993}。在 Garber 看来，体系的崩溃是重大经济事件里“最准确、最普遍地被预言的”之一\cite{garber1993}。
% 据: ei_consensus, ga_predicted

第一个争论，就是缺陷在设计还是在美国的政策。照 Bordo 的结论，一个背着金汇兑本位与可调整钉住这两个致命缺陷的体系，是被世界通胀的加快触发而崩溃的\cite{bordo1993}。而世界通胀的加快，很大程度上是美国通胀先加快的结果\cite{bordo1993}。国际清算银行 1968 年也提到，有人认为大西洋两岸都需要改变财政与货币政策的搭配，可是“政治上的困难似乎是一个障碍”\cite{bis36}。
% 据: bo_verdict, b68_mix

两种解释也有对不上的地方。1968 年以前，美国的通胀率还低于除美国以外七国集团按国民生产总值加权的通胀率\cite{bordo1993}。照 Garber 的分析，难处在识别。美国说逆差来自欧洲资本市场不发达、对流动储备的需求增加；欧洲各国政府说逆差来自美国过度的货币扩张\cite{garber1993}。两种原因在数据里留下的是同样的美元积累，这个识别问题构成了 1960 年代货币争论的“政治层面”\cite{garber1993}。
% 据: bo_us_lower, ga_claim_us, ga_claim_eu, ga_identification

两种诊断开出的药方，后果正好相反。照 Garber 的分析，若逆差来自欧洲对流动储备的需要，互换与借款总安排这些新的信贷额度减少了对储备的需要，会缓和美国的逆差\cite{garber1993}。代价是物价上升，美元兑换黄金也受到威胁\cite{garber1993}。若逆差来自美国过度的货币扩张，它们减少了人们对美元的需要，反而会加速危机\cite{garber1993}。
% 据: ga_alleviate, ga_accelerate

第二个争论是，美元的地位对美国是特权，还是对世界的服务。照 Garber 的概括，法国的批评是，这个体系给了美国作为中央银行的巨大好处\cite{garber1993}。外国人买进它的流动负债或黄金时，并不迫使它采取任何纠正措施，结果“外国的资产被便宜地买走了”\cite{garber1993}。照 Bordo 的说法，作为储备货币国，美国不必为国际收支调整国内经济，联邦储备照例冲销美元的外流，这种不对称招来了不满\cite{bordo1993}。西德则把美国看作在用逆差向顺差国输出通胀\cite{bordo1993}。
% 据: ga_french, ga_cheap, bo_asym, bo_exporting

另一方把美国看作世界的银行。照 Bordo 的概括，在 Despres、Kindleberger 与 Salant（1966）看来，美国在为世界其他地方提供金融中介的服务\cite{bordo1993}。欧洲人向美国借长期资本，因为美国的资本市场更深、利率更低；又把短期存款放在美国的银行，因为收益更高\cite{bordo1993}。这一方的根据是经常账户。除 1959 年以外，美国到 1970 年一直有经常账户的顺差，国际收支的逆差来自资本输出超过了这个顺差\cite{bordo1993}。照 Garber 的说法，这样的逆差显示的只是“资产的重新包装，这是任何金融中介的生意”，美国作为一个整体是在为外国人经营一家银行\cite{garber1993}。
% 据: bo_dks, bo_dks2, bo_current, ga_repackage, ga_bank

两边也各有软肋。法国人说的铸币税，即美国靠别国持有美元得到的好处，照 Bordo 的说法，大部分又以美元存款的利息付给了欧洲\cite{bordo1993}。世界的银行则怕挤兑，美国的货币当局自己也把逆差看作问题，是“因为它对信心的影响”：外国官方的美元越积越多，换成黄金的可能就越大\cite{bordo1993}。银行的类比出自 Kindleberger（1965），特里芬（1966）写文章反驳过\cite{garber1993}。
% 据: bo_fn58, bo_confidence, ga_bank_fn

两边的利益都摆在明处。美国要在海外驻军，要低的长期利率与国内的增长，又不愿改变金价\cite{bis36}，美元的地位让它不必在这几样之间取舍；法国手里的黄金越多，金价翻一番的收益就越大（本书的解读）。照 Bordo 的说法，法国人怨恨美国的金融支配地位，以及他们认为美国得到的铸币税\cite{bordo1993}。两边都有道理，也都在替自己说话。
% 据: b68_limits, bo_french

第三个争论回到\chref{ch:warfinance}的问题。那一章讲过，凯恩斯要债权国分担，最后的协定对顺差国只有稀缺货币条款。照 Bordo 转述的特里芬（1957）的看法，基金组织 1946 年没有宣布美元为稀缺货币是一个错误，否则 1947 年英镑可兑换的灾难本可避免\cite{bordo1993}。西德的马克 1961 年升值 5\%，1969 年升值 9.3\%\cite{bis36}。照 Bordo 的说法，顺差国不愿用升值来调整，越来越认为该由美国调整\cite{bordo1993}。1949—1967 年，十国集团国家的平价很少变动\cite{bordo1993}。
% 据: bo_scarce, b61_reval, b70_de_reval, bo_reluctant, bo_few

升值也不是替逆差国调整。照国际清算银行 1961 年的说法，西德与荷兰的升值，是为了配合稳定国内局势的货币政策；它们主要不是、也不会成为货币受压的国家该采取的措施的“替代”\cite{bis36}。\figref{fig:brettonwoods_parities}里，英镑与法郎降低平价，马克两次升值，日元不动。这段记录可以读作凯恩斯说对了，顺差国不肯动；也可以读作逆差的美国同样没有动。
% 据: b61_substitute, bt_uk, bt_fr, bt_de, bt_jp

\fig{brettonwoods_parities}{四种货币的美元价值，1949 年 10 月—1971 年 8 月（各自 1949 年 10 月的平价为 100）。英镑与法郎降低平价，马克两次升值，日元不变；马克 1969 年 9 月 29 日到 10 月 26 日浮动，1971 年 5 月 9 日起又浮动。数据：Bordo（1993）表 1.2。}{fig:brettonwoods_parities}

\begin{deeper}[史料说明]
\textbf{版本。}Bordo、Garber 与 Eichengreen 的三章，据美国国家经济研究局网站上 Bordo 与 Eichengreen 合编的《布雷顿森林体系回顾》（1993）各章的 PDF。Bordo 的表 1.2 与 Garber 的表 9.1，文字层不保留表格的栏，据页面图像读出。国际清算银行的年报据其网站的扫描本，文字层有“T h e”这样的字距，以及把 1st 读作“ist”、10 读作“io”、1960 读作“i960”的识别错误。Bordo 一章的文字层把 World War II 读作“World War 11”，把 President 读作“F'resident”，把 Triffin 读作“Tiiffin”，注 40 把 scarce 印作“scare”。美国《统计摘要》1964、1968 与 1972 年版的表，据页面图像读出。戴高乐 1965 年 2 月 4 日的话，据美国国务院外交文件第 8 卷第 36 号文件的编者按。

\textbf{不一致之处。}英镑可兑换的中止，Bordo 写作 1947 年 8 月 20 日，国际清算银行写作自 8 月 21 日起，正文用前者。欧洲支付同盟，Bordo 说 1950 年 9 月 19 日成立，国际清算银行说 7 月开始运转，正文用后者。1958 年 12 月底恢复可兑换的国家，Bordo 说 8 国，Garber 的表说 14 个欧洲国家，国际清算银行列举的是同盟成员除希腊、冰岛、土耳其以外再加芬兰。黄金总库，Bordo 的表写作 1961 年 11 月，Garber 的表写作 10 月。美国的份额，Garber 的 50\% 是约定的份额，Bordo 引 Schwartz（1989）的 60\% 是实际卖出的份额，两者不算矛盾。1960 年秋天的金价，Bordo 把 35.20 美元写作美国财政部的买价；国际清算银行记下的是伦敦 9 月中到 10 月中 35.20 至 35.25 美元的报价，美国财政部的卖价则是 35.0875 美元，正文从后者。约定不在高价买金的，Bordo 写作十国集团的货币当局，国际清算银行写作每月在该行开会的各国中央银行，十国的借款总安排 1961 年 12 月才议定，正文从国际清算银行。美国的经常账户，Bordo 说除 1959 年以外到 1970 年都是顺差，Garber 说一直顺差到 1971 年。互换，Bordo 说 1961 年开始，Garber 的表说 1962 年 2 月。吉斯卡尔·德斯坦的说法，Gourinchas 等人系于 1965 年，Lelart 系于 1964 年 9 月，正文两说并存。外国官方的美元负债超过美国的黄金，Bordo 说 1964 年，Garber 说 1960 年代后半期，《统计摘要》在 1967 年底；后者只包括美国银行报告的一年以内的负债，口径可能不同。法国开始换黄金，Garber 说 1962 年第一季度，Bordo 说 1965 年。法郎的贬值，Bordo 的表写作 8 月 10 日，国际清算银行写作 8 月 8 日；马克的新平价，表写作 10 月 26 日，国际清算银行写作 10 月 27 日生效。黄金分成两个市场以后，私人市场金价的高点，Garber 说最高到 42 美元，国际清算银行记下 1969 年 3 月将近 44 美元。《统计摘要》各版对同一年的数偶有修订，1967 年 12 月对外国官方的负债，1968 年版是 140.69 亿美元，1972 年版是 140.27 亿美元，正文与图用后出的一版。

\textbf{转述与间接的证词。}戴高乐的话与财政部的声明，据外交文件编者按的英文概括与引文，正文的中文是从英文译回的。特里芬 1957、1960 与 1966 年的论证，据 Bordo 与 Garber 的转述。肯尼迪 1960 年 10 月 31 日的话，据 Bordo 的注 68。吉斯卡尔·德斯坦与吕夫的说法，据 Gourinchas 等人与 Lelart 两篇后来的论著。英国早些贬值或可免去 1968 年的冲击，是 Johnson（1968）的看法，据 Bordo 的注 52。吕夫（1967）的方案与 Despres、Kindleberger 与 Salant（1966）的看法，据 Bordo 的概括。Bordo、Garber 与 Eichengreen 是同一本书的三章，作者互相引用，不是互相独立的证人。

\textbf{只有一个来源的数。}1946—1949 年欧洲的逆差、美英贷款与法国的多重汇率只据 Bordo；比利时与挪威的数字、1950 年代美国的逆差与黄金只据国际清算银行；法国买走的黄金只据 Garber；特别提款权的分配额只据 Bordo。总库 1967—1968 年卖出的黄金，Bordo 引 Solomon（1976），Garber 的数字相同，两者不是互相独立的证据。美国的黄金与负债只据《统计摘要》各版，它们都出自联邦储备委员会。

\textbf{没有取得的。}基金组织自己的历史（Horsefield 1969，de Vries 1976），imf.org 拒绝访问，没有取得。特里芬 1959 年在国会联合经济委员会的证词与他的书，戴高乐记者招待会的全文（《美国外交政策：现行文件，1965》第 219—221 页），吉斯卡尔·德斯坦的原话，吕夫 1967 年的书，Despres、Kindleberger 与 Salant 的文章，《联邦储备公报》与《银行与货币统计》，也没有取得，所以正文只用转述。1950—1954 年与 1956—1958 年美国的黄金与负债，本书读到的《统计摘要》各版没有，所以图从 1955 年开始，1956—1958 年空缺。稀缺货币条款此后是否被援用过，本书没有找到可以核对的史料。

\textbf{生卒年。}戴高乐、肯尼迪、吉斯卡尔·德斯坦、卡拉汉、克里普斯、凯恩斯与德勃雷据美国国会图书馆的规范档；特里芬与吕夫据德国国家图书馆的规范档（经 lobid.org）。
\end{deeper}
% 据: bo_twothirds, bo_kennedy, bo_tr_rest, bo_scarce, b68_closed, b70_fr_measures, b61_cuts, fr36_text, bo_1947, b48_dates, bo_epu, b59_epu, bo_eight, ga_t91, b59_conv, bt_us, ga_pool50, bo_pool60, bo_oct60, b61_london, b61_selling, bo_g10_3520, b62_measures, b62_gab, bo_current, ga_repackage, bo_fn52, bo_swaps, im_giscard, le_tokyo, bo_1964, ga_wrong, sa72_official, sa72_head, ga_france, bo_convert, bt_fr, b70_fr_deval, bt_de, b70_de_reval, ga_42, b70_gold39, sa68_official, bo_pool_lost, ga_march, bo_dks, bo_rueff, ga_bank_fn

\section{今天还剩什么}\label{sec:brettonwoods-legacy}

\textbf{美元作储备。}第二次世界大战以后，各国把美元当作储备，正是\chref{ch:periphery}讲过的、把储备放在中心国家的老办法。黄金后来退出了，美元仍是各国储备的中心，“过度特权”的争论延续到今天，在\chref{ch:hierarchy}。

\textbf{钉住美元的外围。}\chref{ch:external}讲过，有人把亚洲钉住美元的固定汇率外围看作布雷顿森林体系的再现。这是全球失衡的一种解释，在\chref{ch:moderation}。

\textbf{三选二。}可兑换恢复以后，西德一面钉住汇率，一面想用利率压住通胀，流进来的资本让它两头为难。这是\chref{ch:mundell}的题目。

\textbf{结局。}1971 年 8 月 15 日，美国停止了黄金的兑换\cite{bordo1993}。体系怎样解体，浮动汇率怎样运行，在\chref{ch:floating}。
% 据: bt_us

\begin{takeaways}
\begin{itemize}
\item 布雷顿森林体系里，只有美元按每盎司 35 美元对外国货币当局兑换黄金，其他货币对美元定平价。战后头几年缺的是美元，1947 年英镑恢复可兑换，一个多月就中止了。
\item 1949 年 9 月英镑贬值 30.5\%，降到 2.80 美元，许多货币跟着贬值。代价落在实际工资上；比利时少贬一些，瑞典与挪威用补贴与工资协议分摊，补贴到了上限就涨价。
\item 欧洲支付同盟让西欧多边清算，1958 年底主要货币恢复经常账户的可兑换，1961 年接受第八条的义务。资本开始流动，西德 1961 年 3 月把马克升值 5\%。
\item 特里芬难题：储备要增长，就要美国逆差；美元越多，黄金的覆盖率$c=G^{\mathrm{US}}/D$越低。按《统计摘要》，对外国的负债 1961 年、对官方的负债 1967 年超过了美国的黄金。
\item 美国用黄金总库、互换、罗萨债券、利息平衡税与自愿限制保卫美元，不肯削减驻外军队、改变金价或放弃低利率。1967 年 11 月英镑贬值 14.3\%，1968 年 3 月黄金分成两个市场，特别提款权 1970 年启动。
\item 法国把美元换成黄金，批评美元的“过度特权”；另一方把美国看作世界的银行。缺陷在设计还是在美国的政策，顺差国该不该调整，至今没有定论。
\end{itemize}
\end{takeaways}

\begin{selfcheck}
\begin{enumerate}
\item 在\exref{ex:brettonwoods-triffin}里，若储备的需要每年增长 4\%，其余不变，哪一年$c$跌破 1？
\item 1960 年底，美国的黄金是 178.04 亿美元，对外国官方的负债是 102.12 亿美元。若外国官方把一半的美元换成黄金，美国还剩多少黄金？对官方负债的覆盖率变成多少？
\item 国际清算银行为什么说贬值“等于用间接的办法降低实际工资”？瑞典与挪威的办法有什么不同？
\item 在\exref{ex:brettonwoods-sterling}里，进口的数量不变，出口的数量要增加多少，以英镑计的差额才回到 0？
\item 黄金总库为什么没有挡住美国黄金的流失？
\item 在美元问题上，美国与法国各有什么利益？
\end{enumerate}
\end{selfcheck}

\begin{answers}
\item $D_t=40\times1.04^t-30-0.5t$。第 7 年$D_7\approx52.64-33.5=19.14$，$c_7\approx1.05$；第 8 年$D_8\approx54.74-34=20.74$，$c_8\approx0.96$，第 8 年跌破 1，比 5\% 时晚两年。
\item 换走$102.12/2=51.06$亿美元，黄金剩$178.04-51.06=126.98$亿美元；官方负债剩 51.06 亿美元，覆盖率约 2.49。兑换让两头都减少，覆盖率反而上升，可美国的黄金少了近三成。
\item 进口的东西变贵，工资却不变，同样的工资买到的东西少了。瑞典与挪威用预算补贴压住生活费用，换工会接受工资不变；补贴有上限，挪威 1950 年 4 月只好让黄油、人造黄油与面包涨价。
\item 出口要收 116.7 英镑，数量要增加约 16.7\%，与一美元贵了的幅度相同。
\item 另外 7 家中央银行在总库以外把美元换成黄金，补回自己卖出的黄金；1967 年 12 月到 1968 年 3 月，总库卖掉 30 亿美元的黄金，美国出了 22 亿美元。1968 年 3 月的公报才说，不再把黄金卖给货币当局去补私人市场上卖掉的黄金。
\item 美国可以在海外驻军、维持低利率与国内的增长，又不必改变金价；外国人持有美元，它就不必马上调整。法国持有黄金，金价翻一番对它有利，它也借此反对美国的支配地位。
\end{answers}

1965 年要求回到黄金的戴高乐，1969 年 4 月辞职，法郎的压力更大了\cite{bis36}。1968 年过度扩张的政策正在收回，可它的效果加上工资与其他收入的大幅增加，到 1969 年年中需求仍很旺\cite{bis36}。尽管有外汇管制，国内的储蓄还是借资本外逃与买金流走，上半年国际收支的逆差达 17 亿美元\cite{bis36}。8 月 8 日，法国政府把法郎贬值 11.1\%，有基金组织 10 亿美元的备用安排支持\cite{bis36}。配套的是再削减预算支出，加快征收公司税，提前取消 1968 年推出的 10\% 的投资税收抵免\cite{bis36}。法兰西银行的贴现率 6 月是 7\%，10 月提到 8\%\cite{bis36}。
% 据: b70_degaulle, b70_fr_wages, b70_fr_flight, b70_fr_deval, b70_fr_measures, b70_fr_rate

钱转向了西德。4 月底 5 月初，流入西德的资金略多于 40 亿美元，5 月 9 日政府宣布马克的平价不变，才止住\cite{bis36}。9 月 25 日，为了避免“黑色星期五”，法兰克福市场关闭\cite{bis36}。9 月 29 日到 10 月 26 日，联邦银行停止在干预上限买进美元，马克的汇率浮动，马上升了约 5\%\cite{bis36,bordo1993}。10 月 27 日，新的平价生效，马克升值 9.3\%\cite{bis36}。法国有外汇管制，资本照样外逃；西德钉住汇率，又挡不住流进来的钱。固定汇率、资本自由流动、独立的货币政策，为什么只能三选二？哪些地区适合共用一种货币？
% 据: b70_de_inflow, b70_may9, b70_blackfriday, b70_de_float, b70_at_once, bt_de_c, b70_de_reval
